\documentclass[a4paper]{article}
\usepackage{amsfonts}
\usepackage{amsmath}
\usepackage{amssymb}
\usepackage{graphicx}     
\newcommand{\dint}{\displaystyle\int}

\begin{document}
  
\noindent \large Auteur: Michael Livchits \\
\noindent \large Studierichting: Informatica\\
\noindent \large Oefeningenreeks 5A nr. 5.4\\

\medskip

\normalsize

$I = \displaystyle\dint_{-\tfrac{\pi}{2}}^0 \left(5\cos x-3\sin x\right)dx$\\

Eerst vinden we de onbepaalde integraal\\

$\displaystyle\int \left(5\cos x-3\sin x\right)dx$\\

$ = 5\sin x + 3 \cos x + C$\\

Nu vullen we de grenzen in\\

$I = \left(5\sin 0 + 3 \cos 0\right) - \left(5\sin\left(-\dfrac{\pi}{2}\right)+3\cos\left(-\dfrac{\pi}{2}\right)\right)$\\

$= 3+5 = 8$\\

\begin{thebibliography}{999}
%\bibitem{a} Referentie
\end{thebibliography}
\end{document} 